% auto-generated by scripts/render-tables.py
\documentclass[]{polymathic}
\usepackage{amssymb}
\usepackage{multirow}
\usepackage{bigdelim}
\usepackage{todonotes}
\usepackage{longtable}
\usepackage{tabularray}
\usepackage{subcaption}
\usepackage{wrapfig}
\usepackage[most]{tcolorbox} 
\usepackage{xcolor}
\usepackage{url}
\usepackage{wrapfig}
\usepackage{graphicx}
\usepackage{booktabs}   % nice horizontal rules
\usepackage{tabularx}   % auto-stretch column
\usepackage{caption}    % lets us use \captionof inside 
\renewcommand{\arraystretch}{1.1}  % a touch more row 
\usepackage[sort&compress]{natbib}  % numeric citations like [3, 7–9]
\usepackage{aas_macros}


% \usepackage{draftwatermark}
% \SetWatermarkText{DRAFT}
% \SetWatermarkScale{5.}
% \SetWatermarkColor[gray]{0.85}

\newcommand{\FL}[1]{{\color{violet}FL: #1}}
\newcommand{\JS}[1]{{\color{red}JS: #1}}
\newcommand{\seclead}[1]{{\textbf{\color{polyred}Section Lead: #1}}}

\title{AION-1: Omnimodal Foundation Model for Astronomical Sciences}

\author[*,1,2,3,4]{Liam Parker}
\author[*,5,2]{Francois Lanusse}
\author[*,6]{Jeff Shen}
\author[7]{Ollie Liu}
\author[8]{Tom Hehir}
\author[3]{Leopoldo Sarra}
\author[3]{Lucas Meyer}
\author[9]{Micah Bowles}
\author[2,3]{Sebastian Wagner-Carena}
\author[2]{Helen Qu}
\author[2,3]{Siavash Golkar}
\author[2]{Alberto Bietti}
\author[10]{Hatim Bourfoune}
\author[10]{Nathan Cassereau}
\author[10]{Pierre Cornette}
\author[2,11]{Keiya Hirashima}
\author[2]{Geraud Krawezik}
\author[2]{Ruben Ohana}
\author[3]{Nicholas Lourie}
\author[2,3]{Michael McCabe}
\author[2]{Rudy Morel}
\author[1,8]{Payel Mukhopadhyay}
\author[12]{Mariel Pettee}
\author[2]{Bruno Regaldo-Saint Blancard}
\author[3]{Kyunghyun Cho}
\author[8]{Miles Cranmer}
\author[2,3,6]{Shirley Ho}

\affiliation[1]{University of California, Berkeley}
\affiliation[2]{Flatiron Institute}
\affiliation[3]{New York University}
\affiliation[4]{Lawrence Berkeley National Laboratory}
\affiliation[5]{Université Paris-Saclay,
Université Paris Cité, CEA, CNRS, AIM}
\affiliation[6]{Princeton University}
\affiliation[7]{University of Southern California}
\affiliation[8]{University of Cambridge}
\affiliation[9]{University of Oxford}
\affiliation[10]{IDRIS, CNRS}
\affiliation[11]{RIKEN Center for iTHEMS}
\affiliation[12]{University of Wisconsin--Madison}

\contribution[*]{Equal Contribution.}

\abstract{
While foundation models have shown promise across a variety of fields, astronomy still lacks a unified framework for joint modeling across its highly diverse data modalities. In this paper, we present AION-1, a family of large-scale multimodal foundation models for astronomy. AION-1 integrates heterogeneous imaging, spectroscopic, and scalar data using a two-stage architecture: modality-specific tokenization followed by transformer-based masked modeling of cross-modal token sequences. The model is pretrained on five large-scale surveys: Legacy Survey, Hyper Suprime-Cam (HSC), Sloan Digital Sky Survey (SDSS), Dark Energy Spectroscopic Instrument (DESI), and Gaia. These span more than 200 million observations of stars, galaxies, and quasars. With a single frozen encoder, AION-1 achieves strong results on a broad suite of downstream tasks, including galaxy and stellar property estimation, galaxy morphology classification, similarity-based retrieval, galaxy image segmentation, and spectral super-resolution. We release AION-1 model variants ranging from 300 M to 3.1 B parameters. Beyond astronomy, AION-1 provides a scalable blueprint for multimodal scientific foundation models that can seamlessly integrate noisy, instrument-specific observations. All code, tokenizers, pretrained weights, and a lightweight evaluation suite are released under an open-source license.}

\date{\today}
\titlecorrespondence{Liam Parker: \email{lhparker@berkeley.edu}; Francois Lanusse: \email{francois.lanusse@cnrs.fr}}
\titlecode{\url{https://github.com/PolymathicAI/AION/}}




\usepackage[active,tightpage]{preview}
\PreviewEnvironment{table}
\PreviewEnvironment{table*}
\PreviewEnvironment{subtable}
\PreviewEnvironment{subfigure}
\PreviewEnvironment{wraptable}
\PreviewEnvironment{longtable}
\begin{document}
\begin{preview}
  \centering
  \small
  \begin{tabular}{lcccc}
    \toprule
                       & $\mathbf{T_{\!eff}}$ & $\mathbf{\log g}$ & $\mathbf{[Fe/H]}$ & $\mathbf{v_{\mathrm{mic}}}$ \\
    \midrule
    \multicolumn{5}{@{}l}{AION-1-B}\\[-3pt]
    \hspace{1mm}\textit{Ph}                       & 0.94 & 0.95 & 0.58 & 0.86 \\
    \hspace{1mm}\textit{Ph+Plx+RA/Dec}            & 0.94 & 0.95 & 0.70 & 0.87 \\
    \hspace{1mm}\textit{XP+Ph+Plx+RA/Dec}         & 0.96 & 0.98 & 0.91 & 0.89 \\
    \hspace{1mm}\textit{Sp+Plx}                   & 0.99 & 0.98 & 0.94 & 0.89 \\[2pt]

    \multicolumn{5}{@{}l}{AION-1-L}\\[-3pt]
    \hspace{1mm}\textit{Ph}                       & 0.95 & 0.96 & 0.58 & 0.87 \\
    \hspace{1mm}\textit{Ph+Plx+RA/Dec}            & 0.95 & 0.96 & 0.71 & 0.88 \\
    \hspace{1mm}\textit{XP+Ph+Plx+RA/Dec}         & 0.97 & 0.98 & 0.92 & 0.89 \\
    \hspace{1mm}\textit{Sp+Plx}                   & 0.99 & 0.98 & 0.94 & 0.89 \\[2pt]

    \multicolumn{5}{@{}l}{AION-1-XL}\\[-3pt]
    \hspace{1mm}\textit{Ph}                       & 0.92 & 0.94 & 0.56 & 0.85 \\
    \hspace{1mm}\textit{Ph+Plx+RA/Dec}            & 0.93 & 0.95 & 0.68 & 0.87 \\
    \hspace{1mm}\textit{XP+Ph+Plx+RA/Dec}         & 0.97 & 0.97 & 0.91 & 0.88 \\
    \hspace{1mm}\textit{Sp+Plx}                   & 0.98 & 0.98 & 0.92 & 0.89 \\
    \midrule
    \multicolumn{5}{@{}l}{XGB Baseline}\\[-3pt]
    \hspace{1mm}\textit{Ph}$^{1}$                 & 0.94 & 0.95 & 0.59 & 0.87 \\
    \hspace{1mm}\textit{XP+Ph+Sp}$^{1}$           & 0.99 & 0.98 & 0.89 & 0.89 \\[2pt]

    \multicolumn{5}{@{}l}{ConvNeXT Baseline$^{2}$}\\[-3pt]
    \hspace{1mm}\textit{Sp}                       & 0.99 & 0.98 & 0.95 & 0.89 \\
    \bottomrule
  \end{tabular}
  \captionof{table}{\textbf{$R^2$ ($\uparrow$) for stellar-label prediction.}  
           Inputs: Gaia photometry (\textit{Ph}), low-resolution Gaia XP spectra (\textit{XP}),  
           parallax (\textit{Plx}), celestial coordinates (\textit{RA/Dec}), and high-resolution DESI spectra (\textit{Sp}).  
           $^{1}$Gradient-boosted trees (XGBoost).  
           $^{2}$Baseline convolution–attention network trained on spectra only.}
  \label{tab:stellar_properties_full}
\end{preview}
\end{document}
